% Part II opener: the scope of Chapters 10--16, architecture and implementation.
\noindent Part~II is the applied, code-first half of the book: an engineering
treatise on implementing neural operators for scientific machine learning and
physical simulation. It occupies the ground where the functional analysis of
Part~I meets graphics processing unit (GPU) computing, and its test is concrete
--- whether each discretization decision can be defended against the continuum
statement it is meant to realize. End-to-end implementation patterns are
established in PyTorch and against the official \texttt{neuraloperator}
library, and the architectures are built from source before any of them is
imported: DeepONet, a Deep Operator Network whose branch and trunk networks are
sized, paired, and contracted by hand; the Fourier Neural Operator (FNO), whose
one-, two-, and three-dimensional spectral convolutions are constructed from
real-to-complex fast Fourier transforms (FFTs) and batched tensor contractions;
and the graph neural operator (GNO) together with its geometry-informed (GINO)
composition, for irregular geometries. Each implementation is complete and
annotated, with the commentary naming the tensor shape at every stage, what the
stage costs in arithmetic and memory, and how it fails when the shape is
wrong. Part~II then turns to the decisions that surround a trained model rather
than its layers: discretization-invariant loss functions, infinite-dimensional
data normalization, the stabilization of autoregressive rollout, the numerical
caveats that automatic mixed precision (AMP) introduces, and production
deployment for high-throughput inference through TorchScript and TensorRT. The
chapter sequence follows a training run --- layout, architecture, benchmark,
schedule, serving --- and what it produces is a surrogate whose continuum
claims Part~III puts against physical and biomedical data.
